% =====================================================================
% IV-D: đánh giá lời giải thích (chỉ val); IV-E: MMKG tiếng Việt; IV-F: demo.
% Nguồn số: experiments/xai-metrics-v2/{REPORT.md,summary.json,cases.json}; experiments/xai-metrics-v3/REPORT.md;
% experiments/xai-ground-truth/{REPORT.md,summary.json}; experiments/level3/vnmmkg/xlmr_h14/{a0_metrics,a1_metrics}.json.
\subsection{Kết quả cấp độ 3: đánh giá lời giải thích}
\label{sec:xai}

Mục này kiểm tra lời giải thích $E$ của countonly theo ba câu hỏi: các phần đứng đầu $E$ có thực sự ảnh hưởng tới thứ
hạng không; $E$ của countonly khác $E$ của mô hình gốc \texttt{main\_r3} thế nào; và $E$ đứng ở đâu so với các cách giải
thích khác. Nhắc lại, $E$ là danh sách các phần của đồ thị ảnh hạng 1 (node vật thể hoặc bộ ba), xếp theo đóng góp $t_j$ vào
điểm; để kiểm tra $E$, nhóm bỏ hoặc chỉ giữ các phần đứng đầu rồi chấm lại mô hình (Mục~\ref{sec:xai-metrics}). Mọi số đo
trên val, trung bình $\pm$ std qua 3 seed; giao thức đầy đủ ở Phụ lục~\ref{sec:xai-protocol}.

\begin{table*}[t]
  \caption{Kết quả chính trên val, trung bình $\pm$ std qua 3 seed}
  \label{tab:xai-main}
  \centering
  \footnotesize
  \begin{tabular*}{\textwidth}{@{\extracolsep{\fill}}lcccc@{}}
    \toprule
    & \multicolumn{2}{c}{countonly (v2)} & \multicolumn{2}{c}{\texttt{main\_r3}, mốc so sánh (v2)} \\
    \cmidrule{2-3}\cmidrule{4-5}
    Độ đo & E & NN & E & NN \\
    \midrule
    \multicolumn{5}{@{}l}{\emph{Độ trung thực (fidelity): bỏ $k$ phần đầu của $E$ (Fid$^{+}$), \% câu ảnh mất hạng 1 (cao là tốt)}} \\
    \quad $k=1$ & \pmstd{34,4}{4,8} & \pmstd{4,4}{0,5} & \pmstd{20,5}{1,8} & \pmstd{3,0}{0,2} \\
    \quad $k=3$ & \pmstd{58,1}{4,1} & \pmstd{9,4}{0,9} & \pmstd{34,9}{2,8} & \pmstd{6,5}{0,3} \\
    \quad $k=5$ & \pmstd{68,2}{3,0} & \pmstd{12,9}{0,7} & \pmstd{43,6}{3,1} & \pmstd{9,8}{0,5} \\
    \quad Tính cộng $R_3$: điểm mất / tổng đóng góp của 3 phần đầu & \pmstd{0,189}{0,045} &  & \pmstd{0,063}{0,016} &  \\
    \multicolumn{5}{@{}l}{\emph{Độ trung thực (fidelity): chỉ giữ $k$ phần đầu của $E$ (Fid$^{-}$), hạng của ảnh (thấp là tốt)}} \\
    \quad $k=1$ & \pmstd{4,19}{0,84} & \pmstd{10,16}{2,27} & \pmstd{3,00}{0,03} & \pmstd{9,19}{0,55} \\
    \quad $k=3$ & \pmstd{2,22}{0,74} & \pmstd{8,15}{1,41} & \pmstd{2,16}{0,08} & \pmstd{7,38}{0,36} \\
    \quad $k=5$ & \pmstd{1,86}{0,54} & \pmstd{6,95}{0,96} & \pmstd{2,17}{0,06} & \pmstd{6,89}{0,33} \\
    \multicolumn{5}{@{}l}{\emph{Tính hợp lệ (validity)}} \\
    \quad Câu có $k_{\text{flip}}\le 10$ (\%); trung vị $k_{\text{flip}}$ & \pmstd{80,2}{1,5}; 2 &  & \pmstd{60,3}{2,9}; 3 &  \\
    \quad Vật thể khớp strict (\%) & \pmstd{63,8}{0,9} &  & \pmstd{58,3}{0,3} &  \\
    \quad Vật thể khớp lenient (\%) & \pmstd{70,0}{1,0} &  & \pmstd{63,8}{0,2} &  \\
    \quad Bộ ba khớp đủ chủ thể, đối tượng, vị từ (\%) & \pmstd{2,9}{0,6} &  & \pmstd{5,4}{0,2} &  \\
    \multicolumn{5}{@{}l}{\emph{Độ gọn (sparsity), đơn vị: phần của ảnh}} \\
    \quad $k_{80}$; \% câu đạt trong $k\le 10$ & \pmstd{2,78}{0,27}; 94,5 &  & \pmstd{2,65}{0,14}; 55,7 &  \\
    \bottomrule
  \end{tabular*}
  \par\smallskip
  {\scriptsize\raggedright
    NN: chọn ngẫu nhiên cùng số node vật thể và cùng số bộ ba với $k$ phần đầu của $E$, trung bình 10 lần. Trung vị
    $k_{\text{flip}}$ và trung bình $k_{80}$ chỉ tính trên các câu đạt trong 10 phần. Strict: khớp nhãn hoặc WordNet;
    lenient: tiêu chí rộng hơn của bản đo v1.\par}
\end{table*}

\subsubsection{Kết quả chính}\label{sec:xai-results}
\textbf{Các phần đứng đầu $E$ ảnh hưởng mạnh tới thứ hạng.} Theo Bảng~\ref{tab:xai-main}, bỏ 3 phần đứng đầu $E$ làm ảnh
mất hạng 1 ở 58,1\% câu, gấp sáu lần khi bỏ 3 phần chọn ngẫu nhiên (9,4\%). Ngược lại, chỉ giữ 3 phần đó thì ảnh vẫn ở
hạng trung bình 2,22, so với 8,15 khi giữ 3 phần ngẫu nhiên. Ở 94,5\% câu, chỉ giữ tối đa 10 phần đầu đã còn ít nhất 80\%
phần điểm mà đồ thị mang lại; trên các câu này, trung bình cần 2,78 phần.

\textbf{$E$ của countonly có ảnh hưởng rõ hơn $E$ của \texttt{main\_r3}.} Bỏ 3 phần đứng đầu làm ảnh mất hạng 1 ở 58,1\% câu
với countonly, so với 34,9\% với \texttt{main\_r3}. Tỉ lệ câu tìm được phản thực tế trong 10 phần đầu là 80,2\% so với
60,3\%, và tỉ lệ câu giữ được 80\% $\Delta$ trong 10 phần là 94,5\% so với 55,7\%. Chiều ngược lại không rõ bằng: khi chỉ
giữ 1 phần, ảnh của \texttt{main\_r3} đứng cao hơn (hạng 3,00 so với 4,19). Hai mô hình khác nhau ở bốn thay đổi
(Mục~\ref{sec:method-main}), nên kết quả này phù hợp với lựa chọn trộn điểm thô nhưng chưa tách được tác dụng riêng của nó.

\textbf{$E$ giúp kiểm tra các phần đã đẩy ảnh sai lên đầu.} Khi ảnh hạng 1 sai, bỏ tối đa 10 phần đầu của $E$ làm ảnh này
mất hạng 1 ở 91,8\% câu, so với 68,0\% khi ảnh hạng 1 đúng. Trong các câu có ảnh hạng 1 sai mà ảnh đúng vẫn nằm trong 50
ứng viên (trung bình 5\,069 câu mỗi seed), bỏ 5 phần đầu của $E$ làm ảnh đúng tăng hạng ở 51,9\% câu, và đưa nó lên hạng 1
ở 20,6\%.

\textbf{Ghép vật thể thường đúng loại, ghép bộ ba hiếm khi khớp đủ.} Xét mọi cặp ghép giữa một vật thể của câu và node ảnh
có attention lớn nhất với nó, 63,8\% cặp cùng loại theo tiêu chí strict, và 70,0\% theo tiêu chí lenient rộng hơn. Với các
cặp ghép bộ ba, chỉ 2,9\% khớp đủ cả chủ thể, đối tượng và vị từ.

\subsubsection{Hạn chế của $E$}\label{sec:xai-findings}
Các kết quả trên cho thấy thứ tự của $E$ phản ánh phần nào giữ ảnh ở hạng 1. Hai kết quả sau cho thấy giới hạn của nó.

\textbf{Bộ ba đứng đầu $E$ ít ảnh hưởng hơn node đứng đầu.} Khi phần đứng đầu $E$ là một node vật thể, bỏ nó làm ảnh mất hạng
1 ở \pmstd{46,5}{0,8}\% câu; khi đó là một bộ ba, tỉ lệ chỉ là \pmstd{19,7}{3,2}\% (bỏ ngẫu nhiên: khoảng 4\%). Dù vậy, bộ
ba đứng đầu $E$ ở 43,6\% câu. Với \texttt{main\_r3}, bộ ba đứng đầu ở 70,5\% câu và khoảng cách còn lớn hơn (49,3\% so với
8,4\%). Hai phép bỏ khác nhau về phạm vi: bỏ một node kéo theo mọi bộ ba chứa nó, còn bỏ một bộ ba chỉ bỏ một cạnh. Kết quả
này cũng có thể liên quan đến việc đóng góp của kênh bộ ba phân tán hơn: ba bộ ba lớn nhất chỉ chiếm 38,7\% tổng đóng góp
dương của kênh, so với 68,0\% của ba node lớn nhất. Phụ lục~\ref{sec:xai-cases} minh hoạ hai trường hợp này bằng ảnh
(Hình~\ref{fig:xai-cases}).

\textbf{Tổng đóng góp đánh giá quá cao điểm mất khi bỏ phần.} Đóng góp $t_j$ được tính trên đồ thị gốc. Khi bỏ phần rồi chấm
lại, attention dồn sang các phần còn lại và GAT tính lại biểu diễn của các node, nên điểm mất thật khác tổng đóng góp. Trên
toàn val, tổng điểm mất khi bỏ 3 phần đầu chỉ bằng 18,9\% tổng đóng góp của chúng ($R_3={}$\pmstd{0,189}{0,045}), dù hai
đại lượng vẫn tương quan dương giữa các câu (Spearman \pmstd{0,402}{0,149}). Với \texttt{main\_r3}, tỉ số chỉ còn 0,063 và
gần như không tương quan. Chiều ngược lại cũng không cộng tính: chỉ giữ 5 phần đầu cho điểm cao hơn ảnh gốc (116,6\%
$\Delta$). Occlusion, tức chấm lại sau khi bỏ riêng từng phần, cộng tính tốt hơn nhiều ($R_3={}$\pmstd{0,968}{0,008} trên
mẫu so sánh). Vì vậy demo kèm mỗi dòng của $E$ một phép chấm lại khi bỏ phần đó (cột ``bỏ phần này'', Mục~\ref{sec:demo}).

\subsubsection{So với các cách giải thích khác}\label{sec:xai-compare}
Phụ lục~\ref{sec:xai-compare-results} so $E$ với 15 cách giải thích khác (12 cách ở bản đo v2, 3 cách thêm ở v3) trên một
mẫu 1\,000 câu val, và trên ba bộ có đáp án biết trước.

\textbf{Thứ tự của $E$ gần với Shapley.} Tương quan Spearman giữa thứ tự của $E$ và giá trị Shapley là 0,605, cao hơn
occlusion (0,446) và IG (0,335); chỉ hai cách ước lượng chính giá trị Shapley, GraphSVX (0,754) và GStarX (0,653), cao hơn.
$E$ đạt được điều này mà không cần chấm thêm sau lượt xếp hạng.

\textbf{Về độ trung thực, $E$ đứng sau occlusion và nhóm Shapley.} Deletion AUC, tức phần điểm còn lại khi bỏ dần các phần
đầu (thấp là tốt), của $E$ là 0,070, so với $-0{,}007$ của occlusion và $-0{,}243$ của GraphSVX. Insertion AUC, tức phần
điểm đạt được khi chỉ giữ các phần đầu (cao là tốt), của $E$ là 1,262, ngang IG (1,238) nhưng thấp hơn occlusion (1,447).
Trong các phiên bản được điều chỉnh cho bài toán này, GNN-LRP, PGExplainer và SubgraphX khớp Shapley kém hơn $E$ (Spearman
$\le 0{,}08$).

\textbf{Đáp án biết trước.} Trên bộ mô tả vùng của Visual Genome, GEA của $E$, tức mức trùng giữa các phần đầu của thứ tự và
tập phần đáp án, xấp xỉ Shapley (0,595 so với 0,592); bộ này đo mức khớp với chú thích của người, không phải nguyên nhân của
thứ hạng. Trên dữ liệu tổng hợp, phần đứng đầu $E$ luôn thuộc đáp án ở cả 100 ca mỗi seed. Occlusion đứng đầu ở cả ba bộ.

\textbf{Phản thực tế là điểm yếu của $E$.} Trên 200 câu đầu của mẫu, xét các câu mà bỏ không quá 2 phần đã đủ làm ảnh mất
hạng 1: bỏ dần theo $E$ chỉ đạt đúng số phần bỏ nhỏ nhất ở 59,7\% câu, còn occlusion đạt 90,6\%.

% ---------------------------------------------------------------------
\subsection{Kết quả cấp độ 3: MMKG tiếng Việt}
\label{sec:vnmmkg}

\begin{table}[t]
\centering
\caption{A0, zero-shot: pool 331 ảnh, $\alpha=0$ từ MKG-W val. C: chỉ encoder XLM-R H/14; +KG: thêm kênh đếm
cạnh. Tất định. MRR nhân 100.}
\label{tab:vn-a0}
\footnotesize
\setlength{\tabcolsep}{2.5pt}
\begin{tabular}{@{}lrrrrrrr@{}}
\toprule
 & & \multicolumn{2}{c}{R@1} & \multicolumn{2}{c}{R@10} & \multicolumn{2}{c}{MRR} \\
\cmidrule(lr){3-4}\cmidrule(lr){5-6}\cmidrule(l){7-8}
Nhóm câu & $n$ & C & +KG & C & +KG & C & +KG \\
\midrule
Tất cả & 227 & 23,35 & \textbf{33,92} & 72,69 & \textbf{74,01} & 37,18 & \textbf{45,18} \\
Tự động & 197 & 21,32 & \textbf{29,95} & 70,05 & \textbf{72,59} & 35,21 & \textbf{41,61} \\
Tay, trực tiếp & 12 & \textbf{75,00} & 33,33 & \textbf{100,00} & 75,00 & \textbf{81,25} & 48,90 \\
Tay, gián tiếp & 18 & 11,11 & \textbf{77,78} & 83,33 & \textbf{88,89} & 29,41 & \textbf{81,81} \\
\bottomrule
\end{tabular}
\end{table}

\begin{table}[t]
\centering
\caption{A1: $\alpha$ chọn trên val tiếng Việt, test chạy một lần (103 câu tự động, pool 161 ảnh). Hàng cuối chỉ
để đối chiếu. MRR nhân 100.}
\label{tab:vn-a1}
\footnotesize
\begin{tabular}{@{}lrrrr@{}}
\toprule
Test & R@1 & R@5 & R@10 & MRR \\
\midrule
CLIP (XLM-R H/14) & 37,86 & 72,82 & 79,61 & 52,39 \\
+KG, $\alpha=0{,}7$ (chọn trên val) & 47,57 & 75,73 & 84,47 & 60,41 \\
+KG, $\alpha=0$ (của A0) & 50,49 & 70,87 & 81,55 & 60,46 \\
\bottomrule
\end{tabular}
\end{table}

\textbf{Kết quả.} Ở A0 (Bảng~\ref{tab:vn-a0}), kênh KG nâng R@1 từ 23,35 lên 33,92 (+10,6 điểm). Câu gián tiếp đúng ở hạng 1
tăng từ 2 lên 14 trên 18; câu trực tiếp giảm từ 9 xuống 4 trên 12. Linker nối được ít nhất một thực thể ở 84,14\% số câu. Ở
A1 (Bảng~\ref{tab:vn-a1}), R@1 trên val là 36,17 với $\alpha \in [0;\,0{,}35]$, 37,23 với $\alpha \in [0{,}4;\,0{,}75]$ và
29,79 ở $\alpha=1$; hai đoạn phẳng chỉ cách nhau 1 câu trên 94, và MRR chọn $\alpha=0{,}7$. Trên test, R@1 tăng từ 37,86 lên
47,57 (+9,7 điểm). Với $\alpha=0$, R@1 cao hơn (thêm 3 câu) nhưng R@5 và R@10 thấp hơn, MRR gần bằng; $\alpha$ không được
chọn lại theo test.

\textbf{Kênh đếm cạnh áp dụng được trên dữ liệu mới} dù khác ngôn ngữ, khác encoder và đồ thị nhỏ hơn (325 so với 34\,196
triple KGC-train của MKG-W). Kênh không có tham số học, nên chỉ cần một linker theo tên, một KG và một $\alpha$; ở A0, kênh
này nâng R@1 thêm 10,6 điểm.\label{sec:vn-analysis}

\textbf{Truy vấn gián tiếp.} Với câu ``Ngôi chùa ở Hà Nội được xây dựng dưới thời nhà Lý'', Chùa Một Cột (đúng) lên hạng 1
nhờ ba cạnh \texttt{P131}~$\rightarrow$~Hà Nội, \texttt{P2348} (thời kỳ)~$\rightarrow$~nhà Lý, \texttt{P31}~$\rightarrow$~chùa;
vị trí hành chính và thời kỳ không thấy được trong ảnh (Phụ lục~\ref{app:vn-run}). Bốn câu gián tiếp còn sai đều liên quan
tới quan hệ lớp con, mà kênh đếm cạnh trực tiếp (1-hop) không đi qua. Ở ba câu hỏi ``nhà thờ'' (Ninh Bình, Khánh Hòa, Thành
phố Hồ Chí Minh), từ ``nhà thờ'' được nối vào lớp Q16970, còn đáp án là nhà thờ chính tòa (lớp con Q2977), nên đáp án thiếu
cạnh \texttt{P31} này và không được cộng điểm hơn một nhà thờ thường cùng tỉnh. Câu còn lại nhắc ``di sản thế giới''
(Q9259), còn danh hiệu của đáp án là một thực thể khác (Q43113623).

\textbf{Truy vấn trực tiếp ở $\alpha=0$.} Linker nối từ chỉ loại (``nhà thờ'', ``tháp'', ``chùa'') vào thực thể lớp, nên ảnh
có \texttt{P31} tới đúng lớp đó được cộng một cạnh, còn đáp án thường không có (nhà thờ chính tòa thuộc lớp con; chùa có tháp
không thuộc lớp ``tháp''). Với $\alpha=0$, chênh lệch một cạnh quyết định thứ hạng, bất kể độ tương đồng CLIP với các chi
tiết như ``sơn màu hồng'' hay ``hai tháp chuông''. Từ dùng để so sánh cũng bị nối: câu ``Công trình bằng đá xám có nhiều tầng
mái cong như đình chùa, trên đỉnh có cây thánh giá'' đưa Chùa Bút Tháp lên hạng 1 nhờ cạnh \texttt{P31}~$\rightarrow$~chùa,
trong khi đáp án là một nhà thờ chính tòa. Tỉ lệ câu nối được thực thể (84,14\%) tương đương MKG-W test (80,94\%); các lỗi
phân tích ở trên đều do liên kết từ chỉ loại vào sai lớp.

\textbf{$\alpha$ theo loại câu.} Trên val A1 (chỉ có câu tự động), các giá trị $\alpha$ từ 0 đến 0{,}75 chênh nhau 1 câu. Ở
cùng $\alpha=0$, nhóm gián tiếp thêm 12 câu đúng, nhóm trực tiếp mất 5 câu; một $\alpha$ chung cho mọi câu không tối ưu cho
cả hai nhóm.

% ---------------------------------------------------------------------
\subsection{Demo}
\label{sec:demo}

\file{app/demo_app.py} là một giao diện Gradio, chạy được trên GPU 8\,GB. Người dùng nhập truy vấn tiếng Anh, chọn CLIP thuần
hoặc CLIP + đồ thị, rồi xem top-$k$ ảnh. Mỗi ảnh kèm bảng giải thích như Mục~\ref{sec:method-expl}, khung của các vật thể khớp
vẽ trên ảnh, và cột ``bỏ phần này'': chấm lại ảnh bằng chính mô hình sau khi bỏ phần khớp đó khỏi đồ thị ảnh. Câu gõ vào được
tách bằng FACTUAL-T5 và cùng luật số lượng như lúc đóng gói dữ liệu. Phép kiểm tra ở Mục~\ref{sec:checks} bảo đảm demo cho
cùng kết quả với pipeline đánh giá. Video dự phòng theo kịch bản \file{docs/demo-video-script.md}.
